\documentclass[twoside]{article}

\usepackage{aistats2027}
\usepackage[round]{natbib}
\renewcommand{\bibname}{References}
\renewcommand{\bibsection}{\subsubsection*{\bibname}}
\usepackage{amsmath,amssymb}
\usepackage{booktabs}
\usepackage{graphicx}
\usepackage{subcaption}
\usepackage{microtype}
\usepackage{hyperref}
\usepackage{url}

\graphicspath{{figures/}}

\newcommand{\barrier}{\mathcal{B}}
\newcommand{\loss}{\mathcal{L}}

\begin{document}

\twocolumn[
\aistatstitle{Beyond Permutation Symmetry: Symmetry, Redundancy, and Neural-Network Connectivity}

\aistatsauthor{Anonymous Authors}
\aistatsaddress{Anonymous Institution}
]

\begin{abstract}
The single-basin-modulo-permutation conjecture of
\citet{entezari2021role} proposes that independently trained neural networks
can often be connected by a low-loss linear path after permuting their hidden
units.  We test the limits of this explanation.  Exhaustive enumeration of all
hidden-unit permutations in minimal-width XOR networks yields a counterexample
among gradient-trained solutions: pairs of successful endpoints can remain
linearly incompatible under every permutation.  We then consider a larger
exact symmetry class.  Building on the differentiable Sinkhorn re-basing method
of \citet{guerrero2023rebasin}, we augment permutation alignment with
function-preserving positive ReLU scaling.  We also introduce a method-agnostic
refinement that freezes the permutation returned by any alignment method and
subsequently optimizes only the scales.  On XOR, scale-aware alignment achieves
barriers below the exhaustive permutation-only optimum, showing that its gains
cannot be explained by finding better permutations.  It also lowers barriers
relative to permutation-only alignment across all evaluated VGG/CIFAR-10 and
MLP/Fashion-MNIST architectures.  Finally, we study
training-stage connectivity by extending the same-stage analysis of
\citet{ainsworth2023git} to checkpoints from two independent runs at unequal
epochs, $n\neq N$.  The worse-endpoint barrier, defined as the maximum
interpolation loss in excess of the larger endpoint loss, is negligible early
in training and becomes non-negligible only after a distinct training stage.
Together, these results show that permutation alone does not exhaust the
function-preserving transformations relevant to interpolation and identify
when excess interpolation barriers emerge during training.
\end{abstract}

\section{Introduction}

Are independently trained neural networks genuinely in different basins, or
do they only appear separated because their hidden features use different
coordinates?  The question matters both for the geometry of neural-network
loss landscapes and for operations such as weight averaging.  A high-loss
straight line between two parameter vectors is not conclusive: neural-network
parameterizations are non-identifiable, and functionally equivalent
representatives can have very different interpolation behavior.

The single-basin-modulo-permutation conjecture proposes that most stochastic
gradient descent (SGD) solutions can be permuted into linear mode connectivity
(LMC)~\citep{entezari2021role}.  Re-basing methods make this hypothesis
operational and substantially reduce barriers in MLPs and convolutional
networks~\citep{ainsworth2023git,guerrero2023rebasin}.  Yet a residual barrier
after alignment is ambiguous.  An algorithm may have missed a suitable
permutation, the permutation symmetry class may itself be too small, or the
architecture may lack the representational freedom needed for a low-loss path.
Large-network studies cannot cleanly distinguish these explanations because
their permutation groups are combinatorial and alignment is approximate.

We isolate these mechanisms using XOR networks small enough for exhaustive
permutation search, then test the resulting account in VGG networks.  The
minimal-width experiment first removes the search confound: some trained pairs
remain classification-incompatible under every permutation.  We reproduce
this behavior with full-batch gradient descent (GD) and minibatch SGD.  We then
enlarge the exact symmetry class using positive ReLU scaling.  Starting from
the exhaustive best permutation and optimizing scales alone identifies gains
that cannot be attributed to better correspondence search.  Finally, nonlinear
path and width experiments separate representative choice from architectural
redundancy: symmetry moves within an equivalence class, whereas extra units
create additional routes through parameter space.

Our contributions are:
\begin{itemize}
    \item We give an exhaustive counterexample to universal permutation-based
    LMC for trained networks.  At minimal XOR width, $22/28$ successful pairs
    fail classification-LMC after the best permutation under both full-batch GD
    and minibatch SGD.
    \item We introduce scale-aware alignment and show, under an exhaustive
    permutation oracle, that positive scaling improves or preserves the loss
    barrier for all $67$ tested XOR pairs at widths $3$, $5$, and $7$.
    \item We show an empirical transition in nonlinear connectivity between
    widths two and three, distinguishing representational redundancy from
    symmetry alignment.
    \item We show that scaling improves over permutation-only alignment on
    fixed VGG11/13/16/19 pairs trained on CIFAR-10.
\end{itemize}

\section{Problem Formulation and Related Work}
\label{sec:problem}

\paragraph{Connectivity and interpolation barriers.}
For endpoints $\theta_A$ and $\theta_B$, linear interpolation is
$\theta(t)=(1-t)\theta_A+t\theta_B$.  Following prior work, for a scalar
criterion $C$ (loss or classification error) we measure the excess above the
linear interpolation of endpoint values:
\begin{equation}
\begin{aligned}
\barrier(\theta_A,\theta_B;C)
={}&\max_{t\in[0,1]}\big[C(\theta(t))\\
&-(1-t)C(\theta_A)-tC(\theta_B)\big].
\end{aligned}
\label{eq:barrier}
\end{equation}
The maximum is evaluated on a finite grid in our experiments.  Nonlinear mode
connectivity instead asks whether a general curve
$\gamma:[0,1]\rightarrow\Theta$ joins the endpoints at low loss.  A curved
path may avoid a region crossed by the straight segment, so nonlinear
connectivity does not imply LMC~\citep{draxler2018essentially,garipov2018loss}.

\paragraph{Permutation re-basing.}
For a ReLU MLP block $f(x)=W_2\sigma(W_1x)$ and a permutation matrix $P$,
\begin{equation}
(W_2P^\top)\sigma(PW_1x)=W_2\sigma(W_1x).
\label{eq:perm}
\end{equation}
The same action reorders channels in convolutional networks.  Weight matching
finds layerwise permutations that reduce Euclidean distance, activation
matching aligns representations on data, and differentiable methods optimize
soft assignments~\citep{ainsworth2023git,guerrero2023rebasin}.  REPAIR shows
that correcting activation statistics can further improve interpolation after
permutation~\citep{jordan2023repair}.  These successes support a symmetry-aware
view of loss landscapes, but a residual barrier obtained by a heuristic cannot
establish that the symmetry class itself is insufficient.

\paragraph{Why linear connectivity matters.}
A curved connection establishes geometric connectedness, but a straight path
is the relevant object for parameter averaging.  Checkpoint averaging can
improve generalization within one trajectory~\citep{izmailov2018averaging}, and
averaging related fine-tuned models can be effective without ensemble-time
cost~\citep{wortsman2022model}.  Our focus is the geometric prerequisite of a
low-barrier segment, rather than a claim that an aligned midpoint must
outperform its endpoints.

\section{Scale-Aware Alignment}
\label{sec:method}

Permutation is only one part of the exact symmetry group of a ReLU network.
Positive homogeneity, $\sigma(cz)=c\sigma(z)$ for $c>0$, permits the incoming
weights and bias of a unit to be multiplied by $c$ when its outgoing weight is
divided by $c$.  In matrix form, for positive diagonal $D$,
\begin{equation}
(W_2D^{-1})\sigma(DW_1x)=W_2\sigma(W_1x).
\label{eq:scale}
\end{equation}
Permutation and scaling therefore form an exact monomial action.  Scaling does
not change either endpoint function, but it changes the cross-model terms
created inside a linear parameter interpolation and can therefore change the
barrier.  This motivates studying connectivity under a richer symmetry class
before attributing a residual barrier to genuinely distinct basins.

Sinkhorn re-basing replaces discrete permutation search with a differentiable
relaxation~\citep{guerrero2023rebasin}.  For each hidden layer, it exponentiates
an unconstrained score matrix and alternately normalizes its rows and columns,
producing an approximately doubly stochastic soft assignment.  The temperature
controls the sharpness of this assignment.  Differentiability allows the
assignments at all layers to be optimized jointly for the chosen interpolation
objective; the original method computes their gradients efficiently through
implicit differentiation.  We augment this relaxation with positive scales:
\begin{equation}
S_\ell=\operatorname{Sinkhorn}(Z_\ell/\tau),\qquad
D_\ell=\operatorname{diag}(e^{u_\ell}),
\label{eq:monomial}
\end{equation}
where $Z_\ell$ contains unconstrained assignment scores and $u_\ell$ contains
log-scales.  On layer-$\ell$ outputs we apply $S_\ell D_\ell$ and on the next
layer's inputs the compensating $D_\ell^{-1}S_\ell^\top$.  For evaluation, each
$S_\ell$ is projected to a hard assignment.  The resulting transformation is
exactly function preserving.  During soft optimization, the transpose is the
standard differentiable Sinkhorn approximation rather than an exact inverse.

We consider two variants.  \emph{Joint permutation+scale} optimizes $Z_\ell$
and $u_\ell$ together.  \emph{Scale refinement} first obtains and freezes a
hard permutation, then optimizes $u_\ell$ alone.  On XOR, the frozen
permutation can be the exhaustive optimum, making the incremental contribution
of scale identifiable.  A quadratic penalty on $u_\ell$ controls the unbounded
scale gauge.

The optimization objective is experiment-specific.  For XOR we directly
minimize a differentiable estimate of Eq.~\ref{eq:barrier}: the maximum excess
binary cross-entropy over $31$ interpolation points, plus log-scale
regularization.  For VGG, evaluating the full path at every update is costly,
so alignment minimizes calibration loss at the midpoint; reported barriers
are then computed from the test-set interpolation after hard projection.  This
distinction avoids conflating the scalable proxy used on VGG with the barrier
objective used in the controlled experiment.

\section{Experiment I: Permutation Is Not Enough}
\label{sec:xor}

\paragraph{Question.}
When permutation-based LMC fails, did the alignment algorithm fail, or does no
permutation work?  We answer this in the smallest nontrivial setting: a
one-hidden-layer ReLU network of width $h=2$ trained on the four XOR examples.
There are only $2!=2$ hidden-unit permutations, so the symmetry class can be
searched exactly.

\paragraph{Training and evaluation.}
Models are initialized independently and retained only if they reach $100\%$
endpoint accuracy and low binary cross-entropy.  Our primary protocol uses
full-batch GD.  As a robustness check against the interpretation that the
effect is specific to deterministic optimization, we repeat the experiment
with minibatch SGD using batch size two and reshuffling each epoch.  For each
protocol, eight successful endpoints give $28$ pairs.  We enumerate identity
and swap and evaluate classification error along each interpolation.

\paragraph{Result.}
For both training protocols, only $6/28$ pairs preserve $100\%$ accuracy along
the best-permuted segment; the remaining $22/28$ incur interpolation error
under both possible permutations.  Figure~\ref{fig:xor-counterexample} shows a
successful pair and a representative full-batch failure.  In the failure,
both endpoints correctly implement XOR and have similar decision boundaries,
but no hidden-unit relabeling makes their straight interpolation compatible.
Because the search is exhaustive, this result cannot be attributed to weight
matching, Sinkhorn relaxation, or alignment initialization.

The conclusion is deliberately scoped.  These trained endpoints refute a
universal claim that permutation-based LMC follows automatically from SGD or
GD.  They do not contradict the sufficiently-wide, high-probability form of
the single-basin-modulo-permutation conjecture.  They also complement the
hand-constructed counterexample of
\citet{ainsworth2023git}: LMC failure is not confined to endpoints chosen
algebraically outside a standard training procedure.

\begin{figure*}[t]
    \centering
    \begin{subfigure}[t]{0.49\linewidth}
        \centering
        \includegraphics[width=\linewidth]{success.png}
        \caption{Successful pair: a swap makes the endpoints nearly identical.}
    \end{subfigure}
    \hfill
    \begin{subfigure}[t]{0.49\linewidth}
        \centering
        \includegraphics[width=\linewidth]{unsuccesfull_same_boundary.png}
        \caption{Failed pair: similar boundary, but neither permutation removes the barrier.}
    \end{subfigure}
    \caption{Exhaustive permutation alignment for width-$2$ XOR networks.
    Columns in each panel are interpolation points; top and bottom rows show raw
    and best-permuted interpolation.  Both endpoints classify all four XOR
    points correctly.}
    \label{fig:xor-counterexample}
\end{figure*}

\section{Experiment II: A Richer Symmetry Explains More LMC}
\label{sec:scale-results}

\paragraph{Question and setup.}
Does positive scaling improve compatibility once permutation-search error has
been removed?  We train full-batch GD XOR networks at widths $h=3,5,7$, giving
$21$, $36$, and $10$ valid endpoint pairs.  For each pair we enumerate all
$h!$ permutations and select the one with the smallest sampled loss barrier.
We then freeze this oracle permutation and optimize only positive scales.
Final barriers use $61$ interpolation points.  We also report raw interpolation
and practical Sinkhorn baselines.  Additional optimization details appear in
Appendix~\ref{app:details}.

\subsection{Width lowers the oracle barrier}

Increasing width changes the result.  In the initial exhaustive sweep over
$h\in\{2,3,4,5\}$, the number of pairs maintaining $100\%$ accuracy along the
best-permuted interpolation rises steadily; at $h=5$, every evaluated pair
maintains perfect accuracy.  Perfect accuracy does not mean zero loss barrier:
midpoint confidence can still fall relative to the endpoints.

The larger scale-aware study quantifies this effect in loss.  Table
\ref{tab:xor-scale} shows that the oracle permutation barrier falls from $0.357$
at $h=3$ to $0.111$ at $h=5$ and $0.023$ at $h=7$.  Thus, even before scaling,
width supplies matching flexibility.  Practical Sinkhorn permutation search is
weaker than the oracle at all widths, illustrating why heuristic failure alone
cannot diagnose the absence of an alignment.

\begin{table*}[t]
\caption{XOR loss barrier (mean $\pm$ standard deviation).  Oracle+scale first
selects the exhaustive best permutation, then optimizes only positive scales.
The final row counts pairs for which refinement is no worse than the oracle.}
\label{tab:xor-scale}
\centering
\small
\begin{tabular}{lccc}
\toprule
Method & $h=3$ & $h=5$ & $h=7$ \\
\midrule
No alignment & $0.780\pm0.610$ & $0.608\pm0.465$ & $0.331\pm0.184$ \\
Sinkhorn permutation & $0.444\pm0.325$ & $0.189\pm0.162$ & $0.063\pm0.044$ \\
Oracle permutation & $0.357\pm0.309$ & $0.111\pm0.118$ & $0.023\pm0.015$ \\
Joint permutation+scale & $0.428\pm0.380$ & $0.074\pm0.149$ & $0.009\pm0.009$ \\
Oracle+scale & $\mathbf{0.247\pm0.283}$ & $\mathbf{0.025\pm0.044}$ & $\mathbf{0.005\pm0.003}$ \\
\midrule
Oracle+scale $\leq$ oracle & $21/21$ & $36/36$ & $10/10$ \\
\bottomrule
\end{tabular}
\end{table*}

\paragraph{Scale beyond the oracle permutation.}
The decisive comparison in Table~\ref{tab:xor-scale} is between oracle
permutation and oracle+scale.  The hard assignment is identical in the two
conditions, so the improvements can only come from selecting a different
functionally equivalent scaling of the same endpoint.  Refinement reduces the
mean barrier by $31\%$ at $h=3$, $77\%$ at $h=5$, and $78\%$ at $h=7$.  It is
no worse on every one of the $67$ pairs.  This shows that positive scale is not
merely helping a relaxed optimizer discover a better permutation; it enlarges
the attainable set beyond the exhaustive permutation orbit.

Joint optimization is more mixed at the narrowest width.  It improves over the
permutation oracle on $9/21$, $29/36$, and $8/10$ pairs at widths $3$, $5$, and
$7$, respectively.  At $h=3$, its mean barrier ($0.428$) remains above the
oracle-permutation result ($0.357$), whereas at widths five and seven it beats
the oracle mean.  The difference between joint and oracle-refinement results
exposes an optimization gap: a richer symmetry class can contain better
representatives without a from-scratch solver reliably finding them.

Figure~\ref{fig:xor-profiles} shows the averaged interpolation profiles.  All
curves improve with width, but oracle scale refinement (green) is consistently
lowest.  At $h=7$, its mean maximum excess is only $0.005$, compared with
$0.023$ for the best permutation.  Scaling matters most once width has already
made useful correspondences available: it fine-tunes the balance of incoming
and outgoing weights without changing endpoint predictions.

\begin{figure*}[t]
    \centering
    \begin{subfigure}[t]{0.32\linewidth}
        \includegraphics[width=\linewidth]{3h_xor.png}
    \end{subfigure}
    \begin{subfigure}[t]{0.32\linewidth}
        \includegraphics[width=\linewidth]{h5_xor.png}
    \end{subfigure}
    \begin{subfigure}[t]{0.32\linewidth}
        \includegraphics[width=\linewidth]{h7_xor.png}
    \end{subfigure}
    \caption{Mean XOR interpolation loss at widths $3$, $5$, and $7$: raw
    interpolation (gray), Sinkhorn permutation (orange), exhaustive oracle
    permutation (blue), joint permutation+scale (purple), and oracle
    permutation followed by scale refinement (green).  Exact barriers are in
    Table~\ref{tab:xor-scale}.}
    \label{fig:xor-profiles}
\end{figure*}

\section{Experiment III: Width Changes Nonlinear Connectivity}
\label{sec:nonlinear}

\paragraph{Question and setup.}
Residual linear barriers need not indicate disconnected low-loss sets.  To
test whether width affects available routes rather than only alignment, we
optimize quadratic B\'ezier curves and polygonal chains between independently
trained XOR endpoints.  Paths use up to ten trainable interior points and
minimize the maximum loss over $31$ sampled path locations.

\paragraph{Result.}
At $h=2$, extensive searches find low-loss nonlinear paths only between
essentially identical endpoints; additional bends do not recover nontrivial
connections.  At $h=3$, all studied endpoint pairs are connected by low-loss
paths.  Most are found directly; for one difficult pair, concatenating
low-loss paths through a third trained mode establishes connectivity.
Figure~\ref{fig:nonlinear-width} contrasts representative outcomes.

This result is empirical, not a proof that no width-$2$ path exists outside the
tested families.  Nevertheless, the abrupt change after adding one unit
supports a distinction that endpoint alignment alone cannot make: symmetry
selects among equivalent representatives at fixed capacity, while width
changes the representational resources available along a path.

\begin{figure*}[t]
    \centering
    \includegraphics[width=\linewidth]{unsuccesful_10-14_bezier_3bend.png}
    \vspace{-1ex}
    \includegraphics[width=\linewidth]{3hidden_xor_bezier.png}
    \caption{Nonlinear XOR paths.  \emph{Top:} a representative width-$2$
    pair for which optimizing a B\'ezier path does not remove the high-loss
    region.  \emph{Bottom:} at width $3$, a low-loss path connects endpoints
    with different decision boundaries.  Rows visualize predictions along the
    learned path.}
    \label{fig:nonlinear-width}
\end{figure*}

\section{Experiment IV: VGG/CIFAR-10 Validation}
\label{sec:vgg}

\paragraph{Setup.}
We use fixed, independently trained seed-0/seed-1 pairs of VGG11, VGG13,
VGG16, and VGG19.  Alignment uses a deterministic 1000-example calibration
batch.  We compare raw interpolation, permutation-only Sinkhorn, joint
permutation+scale, and hard permutation followed by scale refinement; test-set
barriers are evaluated after hard projection.

\paragraph{Result.}
Figure~\ref{fig:vgg} reports test-loss barriers for the four architectures.
Permutation-only Sinkhorn reduces the raw barriers from roughly $1.88$--$1.98$
to $0.38$, $0.79$, $1.22$, and $1.21$ for VGG11/13/16/19.  Adding scales
improves on permutation-only in every architecture.  The better of joint and
sequential scaling attains $0.109$, $0.429$, $1.035$, and $1.105$,
respectively.  Relative to permutation-only, these are reductions of about
$71\%$, $45\%$, $15\%$, and $9\%$.

Neither optimization order dominates.  Joint permutation+scale is better for
VGG11 and VGG19; freezing the permutation and refining scales is better for
VGG13 and VGG16.  This echoes XOR: the symmetry class contains better
representatives, but joint nonconvex optimization does not always realize the
full gain.  Substantial barriers also remain for deeper VGGs, so permutation
and scale together are not a complete account of LMC.

The VGG results should be read as a transfer check of the controlled XOR
finding.  Exhaustive permutation search is impossible, only one endpoint pair
is used per architecture, and hyperparameters are selected per method.
Nevertheless, the consistent direction of change across four architectures
supports the practical relevance of scaling beyond the toy setting.

\begin{figure*}[t]
    \centering
    \includegraphics[width=0.78\linewidth]{barplot_vggs.png}
    \caption{Maximum test-loss barrier for independently trained VGG models on
    CIFAR-10.  Both scale-aware variants improve on Sinkhorn permutation-only
    for all four fixed endpoint pairs, although joint and sequential
    optimization trade places across architectures.}
    \label{fig:vgg}
\end{figure*}

\section{Discussion and Limitations}

\paragraph{Symmetry and redundancy play different roles.}
The experiments support a two-part account of apparent basin separation.
Permutation removes discrete coordinate mismatch, and positive scaling moves
continuously along the same exact symmetry orbit.  Neither changes the endpoint
function.  Width, by contrast, changes the architecture's representational
freedom and can create low-loss routes unavailable to a narrower network.
Curved paths further relax the requirement that this route be the straight
segment.  These mechanisms interact, but they are not interchangeable: a
better representative cannot manufacture absent capacity, and additional
capacity does not identify the right representative by itself.

The minimal-width experiment also sharpens the interpretation of
\citet{ainsworth2023git}.  Their hand-constructed counterexample shows that
architecture alone does not guarantee permutation-based LMC, and they suggest
that LMC may be an artifact of SGD.  Our result rules out a broad reading under
which standard gradient-based training itself guarantees alignable endpoints:
neither minibatch SGD nor full-batch GD prevents permutation-incompatible
successful solutions.  It does not rule out an SGD bias toward compatibility
in wider or typical modern networks.  Indeed, the rapid fall in the oracle
barrier with width is consistent with the sufficiently-wide, high-probability
form of the single-basin-modulo-permutation
conjecture~\citep{entezari2021role,ainsworth2023git}.

\paragraph{Optimization versus existence.}
The exhaustive protocol is central to this interpretation.  Without it, a high
barrier after Sinkhorn or weight matching could always be ascribed to search
failure.  Conversely, the gap between joint Sinkhorn and oracle refinement
warns against treating a failed continuous relaxation as evidence that a good
monomial representative does not exist.  Future work should distinguish the
expressivity of a transformation class from the solver used to search it.

\paragraph{Implications for alignment studies.}
A residual barrier should be interpreted in stages.  Search error must first
be separated from limitations of the transformation class.  If permutation is
insufficient, richer exact symmetries should be tested before inferring distinct
basins.  Linear and nonlinear paths should then be distinguished, and width
should be treated as a change in representational resources rather than merely
a larger matching problem.

\paragraph{Limitations.}
The exact conclusions come from small one-hidden-layer XOR networks.  The
minibatch and full-batch protocols are robustness replications, not a controlled
estimate of stochasticity: with the same epoch count, batch size two performs
more parameter updates than full-batch GD.  The VGG
study is non-exhaustive, uses a single pair per architecture, and relies on
method-specific hyperparameter selection.  Our nonlinear path searches can
certify found paths but cannot certify nonexistence.  The scale-aware soft
transformation is exactly function preserving only after projection; during
Sinkhorn optimization it is a relaxation.  Finally, we consider ReLU networks,
for which positive homogeneity holds, and do not address normalization layers,
residual architectures, attention symmetries, or approximate functional
equivalences.  Larger multi-seed studies and comparisons with activation repair
are needed before drawing broad architectural conclusions.

\section{Conclusion}

Permutation symmetry is an important but incomplete explanation of neural-
network connectivity.  Exhaustive XOR search shows that trained endpoints can
remain incompatible under every permutation, for both minibatch SGD and
full-batch GD.  Positive scaling improves even the exhaustive best permutation,
while the width-dependent nonlinear experiment shows that exact symmetries do
not subsume representational constraints.  The same scaling effect transfers
to fixed VGG/CIFAR-10 pairs.  The central distinction is therefore simple:
symmetry determines which representatives we compare; redundancy determines
which paths are available.

\subsection*{AI Use Statement}
Generative AI tools were used to brainstorm presentation choices, improve the
organization and readability of draft text, and assist with programming tasks
including debugging, refactoring, plotting, and \LaTeX{} editing.  They were
not used to generate experimental data, select the reported results, or make
the final scientific interpretations.  The research questions, experimental
design, implementation decisions, analysis, and conclusions were determined by
the authors.  All AI-assisted text and code were reviewed and checked before
inclusion.  The authors take responsibility for the final content, claims,
citations, and artifacts.  This statement must be reconciled with the complete
author team's actual use before submission.

\bibliographystyle{apalike}
\bibliography{references}

\section*{Checklist}

For each question, the selected answer is shown in brackets, followed where
useful by a short justification.

\begin{enumerate}
  \item For all models and algorithms presented, check if you include:
  \begin{enumerate}
    \item A clear description of the mathematical setting, assumptions,
    algorithm, and/or model. [Yes; Sections~\ref{sec:problem}--\ref{sec:method}
    and Appendix~\ref{app:details}.]
    \item An analysis of the properties and complexity (time, space, sample
    size) of any algorithm. [No; empirical behavior rather than asymptotic
    complexity is the focus.]
    \item (Optional) Anonymized source code, with specification of all
    dependencies, including external libraries. [No.]
  \end{enumerate}

  \item For any theoretical claim, check if you include:
  \begin{enumerate}
    \item Statements of the full set of assumptions of all theoretical results.
    [Not Applicable; no formal theorem is claimed.]
    \item Complete proofs of all theoretical results. [Not Applicable.]
    \item Clear explanations of any assumptions. [Yes; the positive-homogeneity
    assumption is stated in Section~\ref{sec:method}.]
  \end{enumerate}

  \item For all figures and tables that present empirical results, check if you include:
  \begin{enumerate}
    \item The code, data, and instructions needed to reproduce the main
    experimental results (either in the supplemental material or as a URL).
    [No; experimental settings are included, but an anonymized code release is
    not yet linked.]
    \item All the training details (e.g., data splits, hyperparameters, how they
    were chosen). [Yes; Appendix~\ref{app:details}.]
    \item A clear definition of the specific measure or statistics and error
    bars (e.g., with respect to the random seed after running experiments
    multiple times). [Yes; Eq.~\ref{eq:barrier}, Table~\ref{tab:xor-scale}, and
    the limitations in Section~\ref{sec:vgg}.]
    \item A description of the computing infrastructure used. (e.g., type of
    GPUs, internal cluster, or cloud provider). [No.]
  \end{enumerate}

  \item If you are using existing assets (e.g., code, data, models) or
  curating/releasing new assets, check if you include:
  \begin{enumerate}
    \item Citations of the creator If your work uses existing assets. [Yes.]
    \item The license information of the assets, if applicable. [No.]
    \item New assets either in the supplemental material or as a URL, if
    applicable. [Not Applicable.]
    \item Information about consent from data providers/curators. [Not
    Applicable; only synthetic XOR and the public CIFAR-10 benchmark are used.]
    \item Discussion of sensible content if applicable, e.g., personally
    identifiable information or offensive content. [Not Applicable.]
  \end{enumerate}

  \item If you used crowdsourcing or conducted research with human subjects,
  check if you include:
  \begin{enumerate}
    \item The full text of instructions given to participants and screenshots.
    [Not Applicable.]
    \item Descriptions of potential participant risks, with links to
    Institutional Review Board (IRB) approvals if applicable. [Not Applicable.]
    \item The estimated hourly wage paid to participants and the total amount
    spent on participant compensation. [Not Applicable.]
  \end{enumerate}
\end{enumerate}

\clearpage
\appendix
\onecolumn
\aistatstitle{Beyond Permutation Symmetry: Symmetry, Redundancy, and
Neural-Network Connectivity: Supplementary Materials}

\section{Reproducibility Details}
\label{app:details}

\subsection{XOR training and exhaustive alignment}

The XOR set is $\{(0,0),(0,1),(1,0),(1,1)\}$ with labels
$\{0,1,1,0\}$.  Models have one ReLU hidden layer, biases, and one output logit.
The primary endpoint protocol uses full-batch GD (implemented with the SGD
optimizer) at learning rate $0.05$ for at most $5000$ epochs.
\texttt{ReduceLROnPlateau} uses factor $0.5$, patience $200$, threshold
$10^{-4}$, and minimum learning rate $5\times10^{-4}$.  Training stops when
loss has not improved by $10^{-4}$ for $500$ epochs.  For each retained pair,
all $h!$ permutations are evaluated.  Interpolation selection prioritizes
classification error in the width-$2$ diagnostic and loss barrier in the
scale-aware study.

The width-$2$ robustness replication uses true minibatch SGD with batch size
two and reshuffles the four XOR examples every epoch.  It otherwise uses the
same initialization range, learning-rate schedule, stopping rule, maximum of
$5000$ epochs, accepted seeds, and exhaustive identity/swap evaluation.  Eight
successful endpoints produce $28$ pairs in each protocol.  Matching epoch
counts does not match optimizer-update counts: the minibatch protocol makes two
updates per epoch and is therefore a robustness comparison, not an isolated
causal test of stochastic gradient noise.  The width-$3/5/7$ scale-aware study
uses the primary full-batch-GD endpoints.

For nonlinear paths, both quadratic B\'ezier curves and polygonal chains are
optimized for $1500$ steps at learning rate $0.05$, with no weight decay and
$31$ path samples.  The default path has one trainable interior point; harder
pairs use up to ten.

\subsection{XOR scale-aware optimization}

Final interpolation curves use $61$ evaluation points.  XOR alignment directly
optimizes the maximum sampled excess binary cross-entropy over $31$
interpolation points, rather than midpoint loss alone.  The default
permutation-only Sinkhorn run uses Adam for $500$ steps,
learning rate $0.01$, temperature $1$, $20$ Sinkhorn iterations, and identity
strength $1$.  Scale refinement uses Adam for $500$ steps, learning rate
$0.01$, and log-scale regularization $10^{-3}$.  Joint optimization uses the
same default budget.  If a default run misses its target, a bounded search over
learning rates, temperatures, and regularization is used; settings are selected
per pair.  The scale target is barrier $\epsilon=0.01$, while permutation-only
Sinkhorn targets that pair's exhaustive barrier.

\subsection{VGG experiments}

For the scale-aware comparison, fixed seed-0/seed-1 endpoint pairs are used for
VGG11/13/16/19.  Alignment uses a deterministic batch of $1000$ images, 20
Sinkhorn iterations, and midpoint loss.  Separate bounded sweeps select the
number of optimization steps, learning rate, temperature, and scale
regularization for each architecture and method.  The permutation-then-scale
baseline freezes the projected hard permutation and optimizes only log-scales.

\section{Additional Pairwise Results}

Table~\ref{tab:pairwise} separates the practical optimization gap from the gain
provided by the symmetry class.  Sinkhorn permutation-only exactly matches the
exhaustive oracle on progressively fewer pairs as width grows.  Joint
permutation+scale often beats the permutation oracle, but oracle scale
refinement is the most reliable condition.

\begin{center}
\captionof{table}{Pairwise XOR comparisons.  Lower loss barrier is better.}
\label{tab:pairwise}
\small
\begin{tabular}{lccc}
\toprule
Comparison & $h=3$ & $h=5$ & $h=7$ \\
\midrule
Oracle+scale $\leq$ oracle permutation & $21/21$ & $36/36$ & $10/10$ \\
Joint permutation+scale $\leq$ oracle permutation & $9/21$ & $29/36$ & $8/10$ \\
Sinkhorn permutation $=$ oracle permutation & $9/21$ & $7/36$ & $0/10$ \\
Oracle+scale $\leq$ joint permutation+scale & $15/21$ & $25/36$ & $5/10$ \\
\bottomrule
\end{tabular}
\end{center}

\end{document}
